% Part: first-order-logic
% Chapter: natural-deduction
% Section: soundness-identity

\documentclass[../../../include/open-logic-section]{subfiles}

\begin{document}

\olfileid{fol}{ntd}{sid}

\olsection{Correctheid met \usetoken{S}{identity}}

\begin{prop}
Natuurlijke deductie blijft correct als we regels voor $\eq$
toevoegen.
\end{prop}

\begin{proof}
Elke !!{formula} van de vorm $\eq[t][t]$ is geldig. In elke
!!{structure}~$\Struct M$ geldt namelijk $\Sat{M}{\eq[t][t]}$. We
nemen aan dat de term $t$ gesloten is, dus geen variabelen bevat.
Welke objecten we aan variabelen toekennen, maakt daarom geen verschil.

Stel dat de laatste stap in !!a{derivation} \Elim{\eq} is. Dan ziet de
!!{derivation} er zo uit:
\begin{prooftree}
  \AxiomC{$\Gamma_1$}
  \RightLabel{$\delta_1$}
  \DeduceC{$\eq[t_1][t_2]$}
  \AxiomC{$\Gamma_2$}
  \RightLabel{$\delta_2$}
  \DeduceC{$!A(t_1)$}
  \RightLabel{\Elim{\eq}}
  \BinaryInfC{$!A(t_2)$}
\end{prooftree}
De uitgangsformules $\eq[t_1][t_2]$ en $!A(t_1)$ zijn !!{derive}d uit
de !!{undischarged} aannamen~$\Gamma_1$ en $\Gamma_2$, in die
volgorde.
We willen laten zien dat $!A(t_2)$ volgt uit $\Gamma_1 \cup
\Gamma_2$. Neem
!!a{structure}~$\Struct{M}$ waarvoor $\Sat{M}{\Gamma_1 \cup
  \Gamma_2}$. De inductiehypothese geeft dan
$\Sat{M}{!A(t_1)}$ en $\Sat{M}{\eq[t_1][t_2]}$. Dus
$\Value{t_1}{M} = \Value{t_2}{M}$. Kies een willekeurige toekenning
$s$ van objecten aan variabelen en stel
$m = \Value{t_1}{M} = \Value{t_2}{M}$. Volgens
\olref[fol][syn][ext]{prop:ext-formulas} geldt
$\Sat{M}{!A(t_1)}[s]$ precies dan als
$\Sat{M}{!A(x)}[\Subst{s}{m}{x}]$, en dit geldt precies dan als
$\Sat{M}{!A(t_2)}[s]$. We weten dat $\Sat{M}{!A(t_1)}$. Daarom geldt
ook $\Sat{M}{!A(t_2)}$.
\end{proof}
\end{document}
